\chapter{不完备反馈下的语义交付与能量协同控制}
\label{rq3:chapter}

\section{引言}

在第~\ref{rq1:chapter} 章资源管理与第~\ref{rq2:chapter} 章内容选择的基础上，本章研究消息随机传输、确认不完备条件下的语义交付。图~\ref{rq3:fig:system} 展示了发送端依据本地记录与累计确认，在补发、切换候选和等待之间决策的反馈过程。

\begin{figure}[htbp]
\centering
\begin{tikzpicture}[x=1mm,y=1mm,scale=1.15,transform shape,font=\scriptsize,
  box/.style={draw,rounded corners=1.2mm,align=center,inner sep=3pt,minimum height=11mm},
  flow/.style={-{Stealth[length=1.5mm]},line width=.55pt},
  ack/.style={flow,dashed}]
  \node[box,text width=21mm] (a) at (13,34) {感知节点 A\\局部语义证据};
  \node[box,text width=21mm] (b) at (13,17) {感知节点 B\\局部语义证据};
  \node[box,text width=21mm] (c) at (13,0) {其他感知节点\\局部语义证据};
  \node[box,text width=31mm,minimum height=30mm] (ctrl) at (58,17) {C-BFC 控制器\\[3pt]接收状态概率\\语义效用与能量价格\\发送／重传／等待／停止};
  \node[box,text width=23mm,minimum height=23mm] (rx) at (104,17) {接收中心\\唯一证据融合\\任务分类结果};
  \draw[flow] (a.east) -- ++(6,0) |- ($(ctrl.west)+(0,8mm)$);
  \draw[flow] (b.east) -- (ctrl.west);
  \draw[flow] (c.east) -- ++(6,0) |- ($(ctrl.west)+(0,-8mm)$);
  \draw[flow] (ctrl.east) -- node[above,align=center] {无线数据\\链路} (rx.west);
  \draw[ack] (rx.north) -- ++(0,17) -| node[pos=.25,above] {延迟或丢失的累计确认} (ctrl.north);
  \node[align=center] at (61,-9) {节点发送能量与任务截止时间约束};
\end{tikzpicture}

\bicaption{语义可靠交付的逻辑结构与反馈关系}{Logical architecture and feedback flow of reliable semantic delivery}
\label{rq3:fig:system}
\end{figure}

本章提出一致性感知的信念反馈控制方法 C-BFC（Consistency-aware Belief Feedback Control）。在固定候选与融合规则下，通过效用网络及分类一致性投影估计任务收益，以接收信念描述未确认状态，再用两步前瞻和跨任务虚拟能量队列协调交付收益、任务时限与通信能耗，并分析不完备反馈下的决策条件与性能取舍。

\section{系统模型与问题建模}
\label{rq3:sec:model}

\subsection{网络与语义任务模型}
\label{rq3:subsec:tasks}

考虑由 $N$ 个感知节点与一个接收中心组成的协同通信系统。感知节点集合记为 $\mathcal{N}=\{1,\ldots,N\}$，接收中心记为 $d$。系统依次处理 $N_{\mathrm{task}}$ 个任务，任务索引为 $k=1,\ldots,N_{\mathrm{task}}$。每个任务具有相对截止时间 $D_k$，其通信过程结束后，系统转入下一个任务。节点发送能量在任务之间连续扣减，接收状态和确认记录则在每个任务开始时重新初始化。

对于任务 $k$，上层候选选择模块根据局部观测与语义粒度给出候选证据集合
\begin{equation}
  \mathcal{M}_k=\{m_{k,1},m_{k,2},\ldots,m_{k,K_k}\}.
  \label{rq3:eq:candidates}
\end{equation}
每条证据 $m$ 包含发送节点 $s(m)$、长度 $\ell_m$、语义粒度和接收端信念更新量。不同粒度对应标签、简洁声明、摘要和完整证据。本章固定候选内容、候选次序和融合规则，控制器只调整这些证据的发送过程。选择来自另一节点的候选，即对应在给定候选集合内切换发送节点。

设接收中心实际已经获得的唯一证据集合为 $S_{k,t}\subseteq\mathcal{M}_k$，其中 $t$ 表示当前任务内的决策时刻。接收状态空间为
\begin{equation}
  \states_k=2^{\mathcal{M}_k},\qquad |\states_k|=2^{K_k}.
  \label{rq3:eq:state-space}
\end{equation}
接收中心按固定候选次序融合 $S_{k,t}$ 中的证据，并输出类别 $c_k(S_{k,t})$。相同证据无论成功接收多少次，都只参与一次融合。这一设定区分了通信层的重复发送与语义层的新增信息，使最终分类由证据集合决定，不随同一集合的到达次序变化。

记任务参考类别为 $y_k$，集合 $S$ 对应的真实任务效用定义为
\begin{equation}
  u_k(S)=\ind\{S\ne\varnothing,\ c_k(S)=y_k\}.
  \label{rq3:eq:true-utility}
\end{equation}
空接收集合的效用设为零，表示任务没有通过本次通信获得有效服务。$y_k$ 只用于离线效用学习和实验评价，不属于在线控制器的观测。由于在线无法直接计算式~\eqref{rq3:eq:true-utility}，第~\ref{rq3:subsec:utility} 节通过学习模型估计各接收集合的效用。

本章考虑短任务内持续存在的单跳接触。数据动作串行执行，反馈使用独立控制信道；不同尝试之间的速率扰动相互独立。候选描述及条件融合效果作为预先给定的调度信息，用于计算假设接收集合的语义特征。实验单独计入数据与确认的发送能耗，候选元数据分发、感知推理和飞行能耗不在控制范围内。这些假设使评价集中于给定语义内容后的交付决策。

\subsection{分组传输与随机送达模型}
\label{rq3:subsec:channel}

设候选 $m$ 的有效预测速率为 $\bar r_{k,m}$。链路预测由节点距离、有效带宽和发射功率给出。在本章使用的视距链路近似下，
\begin{equation}
  h_{k,i}=\frac{\rho_0}{\max\{d_{k,i},1\}^2},\qquad
  \bar r_{k,i}=W_k\log_2\left(1+\frac{P^{\mathrm d}h_{k,i}}{N_0W_k}\right),
  \label{rq3:eq:rate}
\end{equation}
其中 $d_{k,i}$ 为节点 $i$ 到中心的距离，$\rho_0$ 为参考信道增益，$N_0$ 为噪声功率谱密度，$P^{\mathrm d}$ 为数据发射功率，$W_k$ 表示任务对应的有效带宽。短任务内以该预测值计算动作时长，跨任务的距离与可用带宽变化则反映在不同的预测速率中。

语义长度以 token 数表示。设平均每个 token 占用 $b$ 字节，分组有效载荷上限为 $L_{\mathrm{MTU}}$ 字节，分组头部长度为 $H$ 字节，则实际发送比特数为
\begin{equation}
  B_m=8\left(b\ell_m+H\left\lceil\frac{b\ell_m}{L_{\mathrm{MTU}}}\right\rceil\right).
  \label{rq3:eq:bits}
\end{equation}
该换算同时保留了消息长度与协议头部开销。一次发送使用的时长为
\begin{equation}
  \tau_{k,m}=\min\left\{\tau_{\max},\ \kappa\frac{B_m}{\bar r_{k,m}}\right\},
  \label{rq3:eq:duration}
\end{equation}
其中 $\kappa$ 是发送时长裕量，$\tau_{\max}$ 是单次占用时长上限。控制器只允许完整时长能够容纳于剩余任务窗口内的发送动作。

为表示每次尝试的速率波动，令实际速率满足
\begin{equation}
  r_{k,m}^{(a)}=\bar r_{k,m}\exp\left(\sigma_r Z_a-\frac{\sigma_r^2}{2}\right),
  \qquad Z_a\sim\mathcal{N}(0,1),
  \label{rq3:eq:random-rate}
\end{equation}
其中 $a$ 为尝试编号，$\sigma_r$ 控制波动强度。若在预留时长内能够传完全部比特，则本次证据送达。相应成功概率为
\begin{equation}
  p_{k,m}=\Prb\{r_{k,m}^{(a)}\tau_{k,m}\ge B_m\}
  =\Phi\left(\frac{\ln(\bar r_{k,m}\tau_{k,m}/B_m)-\sigma_r^2/2}{\sigma_r}\right),
  \label{rq3:eq:success-probability}
\end{equation}
其中 $\Phi(\cdot)$ 为标准正态分布函数；当 $\sigma_r=0$ 时，上式退化为确定的容量判定。实验直接依据式~\eqref{rq3:eq:success-probability} 进行分组成功抽样。相同证据的不同发送尝试独立抽样，失败尝试消耗完整的预留发送时长。

本模型将未完整送达的证据视为擦除，不保存失败传输中的软信息。因此，重传的作用是提高证据至少成功送达一次的概率，而接收端不会对多个失败波形执行 Chase 合并或增量冗余译码。

\subsection{累计确认与可观测历史}
\label{rq3:subsec:feedback}

每次发送尝试结束后，接收中心生成一个累计确认帧。协议假设短控制头能够标识本次调度的发送尝试，因此即使载荷失败，也可返回包含当前接收状态的累计位图。反馈帧携带尝试序号 $j$ 及该时刻的接收集合快照 $S^{\mathrm{ack}}_j$，既可确认已收到的证据，也可表明某条已经尝试的证据仍未送达。

记第 $j$ 次数据发送结束时间为 $t_j^{\mathrm{end}}$，反馈帧长度为 $B^{\mathrm f}$ 比特。反馈按对应链路的预测速率串行发送，完成时间满足
\begin{equation}
  v_j=\max\{t_j^{\mathrm{end}},v_{j-1}\}+\tau_j^{\mathrm f},
  \qquad \tau_j^{\mathrm f}=\frac{B^{\mathrm f}}{\bar r_{k,m_j}},
  \label{rq3:eq:feedback-serialization}
\end{equation}
其中 $v_{-1}=0$。在串行发送之后，反馈以概率 $\epsilon_{\mathrm f}$ 丢失；若未丢失，则经过额外时延 $L_j$ 到达发送端，
\begin{equation}
  L_j\sim\mathcal{U}\left(\frac{\bar d_{\mathrm f}}{2},\frac{3\bar d_{\mathrm f}}{2}\right),
  \qquad t_j^{\mathrm{ack}}=v_j+L_j.
  \label{rq3:eq:feedback-delay}
\end{equation}
数据成功、反馈丢失和反馈额外时延分别抽样。累计确认可能乱序到达，控制器保留序号最大的有效快照，较旧位图不覆盖较新信息。

在决策时刻 $t$，可观测历史 $\history_{k,t}$ 包含候选描述、链路预测、已执行发送记录、已经到达的确认、剩余时间和能量状态。尚未到达的位图内容以及真实接收集合不属于该历史。记
\begin{equation}
  b_{k,t}(S)=\Prb\{S_{k,t}=S\mid\history_{k,t}\},\qquad S\in\states_k,
  \label{rq3:eq:belief-definition}
\end{equation}
为控制器对接收状态的信念。由于反馈过程与数据是否成功独立，确认尚未到达不能直接视为否定确认；它主要改变未来等待能够获得信息的概率。

\subsection{能量状态与动作约束}
\label{rq3:subsec:energy}

节点 $i$ 的初始发送能量为 $E_i^{(0)}$。一次数据发送及其反馈所消耗的能量分别为
\begin{equation}
  e_{k,m}^{\mathrm d}=P^{\mathrm d}\tau_{k,m},\qquad
  e_{k,m}^{\mathrm f}=P^{\mathrm f}\tau_{k,m}^{\mathrm f}.
  \label{rq3:eq:transmit-energy}
\end{equation}
数据发送能量从对应感知节点的剩余电量中扣除。接收中心采用独立供能，反馈发送能量计入系统通信成本。已经生成的反馈帧即使丢失或晚于任务结束才到达，也计入反馈能耗，以反映协议实际发起的发送开销。

设第 $k$ 个任务中节点 $i$ 的数据发送能耗为 $e_i(k)$，则跨任务能量状态满足
\begin{equation}
  E_i(k+1)=E_i(k)-e_i(k),\qquad E_i(1)=E_i^{(0)}.
  \label{rq3:eq:energy-evolution}
\end{equation}
任务内每执行一次发送便立即扣减相应电量，式~\eqref{rq3:eq:energy-evolution} 是任务尺度上的汇总形式。

在时刻 $t$，可选动作包括发送候选 $m$、等待长度 $w$ 以及停止。发送已经尝试过的候选表示重传，发送尚未尝试的候选表示新增证据传输。若当前已执行尝试数为 $n_{k,t}$、尝试数上限为 $A_{\max}$，则发送动作的可行集合为
\begin{equation}
  \begin{split}
  \actions^{\mathrm{send}}_{k,t}=\{m\in\mathcal{M}_k:\;&
  \tau_{k,m}\le D_k-t,\\[-2pt]
  &e_{k,m}^{\mathrm d}\le E_{s(m)}(k,t),\quad n_{k,t}<A_{\max}\}.
  \end{split}
  \label{rq3:eq:action-mask}
\end{equation}
等待动作满足 $0<w\le D_k-t$，停止动作始终可执行。硬约束直接排除无法完成的动作，避免仅通过代价惩罚近似电量和截止期限制。

\subsection{优化目标与求解思路}
\label{rq3:subsec:objective}

记策略 $\pi$ 在任务 $k$ 上的停止时刻为 $T_k$，终止接收集合为 $S_k^{\mathrm{end}}$，该任务总通信能耗为
\begin{equation}
  E_k^{\mathrm{comm}}=\sum_{a\in\mathcal{T}_k}
    \left(e_{k,m_a}^{\mathrm d}+e_{k,m_a}^{\mathrm f}\right),
  \label{rq3:eq:total-communication-energy}
\end{equation}
其中 $\mathcal{T}_k$ 是实际执行的数据发送尝试集合。以正确完成收益与通信消耗之间的权衡为目标，可将有限任务周期内的问题写为
\begin{equation}
  \begin{aligned}
    \max_{\pi}\quad &\frac{1}{N_{\mathrm{task}}}\E_{\pi}\left[
      \sum_{k=1}^{N_{\mathrm{task}}}\left(u_k(S_k^{\mathrm{end}})-\lambda_E
      \frac{E_k^{\mathrm{comm}}}{E_{\mathrm{ref}}}\right)\right] \\
    \text{s.t.}\quad &\sum_{k=1}^{N_{\mathrm{task}}} e_i(k)\le E_i^{(0)},\quad i\in\mathcal{N},\\
    &0\le T_k\le D_k,\quad |\mathcal{T}_k|\le A_{\max},\\
    &a_{k,t}\text{ 仅依赖 }\history_{k,t}\text{，且满足式~\eqref{rq3:eq:action-mask}}.
  \end{aligned}
  \label{rq3:eq:global-problem}
\end{equation}
其中 $E_{\mathrm{ref}}>0$ 为固定能量尺度，$\lambda_E\ge0$ 为系统设计中的能耗偏好。式~\eqref{rq3:eq:global-problem} 描述研究目标，并不意味着能够直接获得其全局最优解：真实任务效用在线不可见，未来数据与反馈具有随机性，节点电量还会将不同任务的决策关联起来。

本章采用分层近似求解。首先利用语义效用网络替代不可观测的即时任务正确性；其次把可见历史压缩为接收信念与资源状态；最后以随资源余量变化的能量价格和有限动作前瞻选择当前动作。这样可以保留式~\eqref{rq3:eq:global-problem} 中的交付收益与资源约束，同时避免枚举完整任务周期的反馈分支。后续算法中的能量权重是局部控制参数，需要通过验证任务选择，不与式~\eqref{rq3:eq:global-problem} 的全局最优乘子等同。

\begin{table}[t]
  \centering\small
  \bicaption{语义交付控制模型的主要符号}{Main notation of the semantic delivery control model}
  \label{rq3:tab:notation}
  \begin{tabularx}{\linewidth}{@{}lY@{}}
    \toprule
    符号 & 含义\\
    \midrule
    $N,\ N_{\mathrm{task}}$ & 感知节点数及连续执行的任务总数\\
    $\mathcal{M}_k,\ K_k$ & 任务候选证据集合及其规模\\
    $S_{k,t},\ b_{k,t}(S)$ & 实际接收集合及其条件概率\\
    $c_k(S),\ u_k(S)$ & 接收集合对应的预测类别与真实任务效用\\
    $\hat q_\theta(S),\ q_k(S)$ & 原始预测效用与一致性修正效用\\
    $\bar r_{k,m},\ p_{k,m}$ & 预测链路速率与单次送达概率\\
    $\tau_{k,m},\ D_k$ & 数据发送时长与任务截止时间\\
    $\bar d_{\mathrm f},\epsilon_{\mathrm f}$ & 反馈平均额外时延与反馈丢失概率\\
    $E_i(k,t),\ Z_i(k)$ & 节点剩余发送能量与虚拟能量队列\\
    $\pi_{k,m}$ & 一次发送及其反馈的归一化资源价格\\
    $\eta_E,\eta_T$ & 局部控制的能量权重与时间权重\\
    \bottomrule
  \end{tabularx}
\end{table}

\section{一致性感知的信念反馈控制算法}
\label{rq3:sec:method}

\subsection{总体设计}

C-BFC 由语义效用估计、接收状态滤波和反馈资源控制三个部分组成，如图~\ref{rq3:fig:controller} 所示。任务开始时，控制器计算候选证据的所有可能接收集合及其语义特征，得到满足分类一致性的效用表；任务执行过程中，根据发送记录和已到达的累计确认更新接收信念，再结合能量价格与动作约束，评价发送、等待和停止。效用表在同一任务内保持不变，接收信念与可执行动作随传输过程更新。

\begin{figure}[t]
  \centering
  \begin{tikzpicture}[x=1mm,y=1mm,
  font=\fontsize{9.5bp}{12bp}\selectfont,
  cbbox/.style={draw,line width=.45pt,align=center,inner sep=2mm},
  cbarrow/.style={-{Stealth[length=1.6mm,width=1.2mm]},line width=.5pt},
  cbgroup/.style={draw,line width=.55pt},
  cbtitle/.style={font=\fontsize{10bp}{13bp}\selectfont\bfseries}]
  \draw[cbgroup] (0,66) rectangle (66,132);
  \draw[cbgroup] (74,66) rectangle (137,132);
  \node[cbtitle] at (33,127) {任务效用表构造};
  \node[cbtitle] at (105.5,127) {在线状态更新};
  \node[cbbox,text width=56mm,minimum height=15mm] (sets) at (33,114)
    {候选证据 $\mathcal M_k$，$K\le4$\\枚举 $2^K$ 个接收集合};
  \node[cbbox,text width=56mm,minimum height=14mm] (network) at (33,95)
    {非空集合特征与效用网络\\$29\rightarrow128\rightarrow96\rightarrow1$};
  \node[cbbox,text width=56mm,minimum height=14mm] (utility) at (33,76)
    {同类别非空集合的效用取均值\\$q_k(S)$，$q_k(\varnothing)=0$};
  \draw[cbarrow] (sets.south) -- (network.north);
  \draw[cbarrow] (network.south) -- (utility.north);
  \node[cbbox,text width=23mm,minimum height=24mm] (history) at (89.5,108)
    {累计确认\\发送记录\\送达率 $p_{k,m}$};
  \node[cbbox,text width=23mm,minimum height=24mm] (resources) at (121.5,108)
    {剩余时限\\节点能量\\$h,E_i,Z_i$};
  \node[cbbox,text width=23mm,minimum height=18mm] (belief) at (89.5,78)
    {接收状态滤波\\$b_{k,t}(S)$};
  \node[cbbox,text width=23mm,minimum height=18mm] (price) at (121.5,78)
    {动作可行性\\发送价格 $\pi_{k,m}$};
  \draw[cbarrow] (history.south) -- (belief.north);
  \draw[cbarrow] (resources.south) -- (price.north);
  \node[cbbox,text width=129mm,minimum height=12mm,fill=black!4] (expected) at (68.5,53)
    {期望任务效用 $U_k(b)=\sum_{S}b(S)q_k(S)$\\在剩余时间、电量和尝试次数约束下评价可行动作};
  \draw[cbarrow] (utility.south) -- (33,59);
  \draw[cbarrow] (belief.south) -- (89.5,59);
  \draw[cbarrow] (price.south) -- (121.5,59);
  \node[cbbox,text width=36mm,minimum height=20mm] (send) at (23,28)
    {两步发送前瞻\\新发或重传\\$V_m^{\mathrm{send}}(b;h)$};
  \node[cbbox,text width=39mm,minimum height=20mm] (wait) at (72.5,28)
    {反馈等待价值\\至多三个等待候选\\$V^{\mathrm{wait}}(w;b,h)$};
  \node[cbbox,text width=28mm,minimum height=20mm] (stop) at (120.5,28)
    {停止当前任务\\净增益为零};
  \draw (68.5,47) -- (68.5,42);
  \draw (23,42) -- (120.5,42);
  \draw[cbarrow] (23,42) -- (send.north);
  \draw[cbarrow] (72.5,42) -- (wait.north);
  \draw[cbarrow] (120.5,42) -- (stop.north);
  \node[cbbox,text width=129mm,minimum height=16mm,fill=black!4] (action) at (68.5,0)
    {选择最大正收益动作并执行；否则停止\\更新发送记录、剩余时间和电量；任务结束后更新 $Z_i(k+1)$};
  \draw (send.south) -- (23,12) -- (120.5,12) -- (stop.south);
  \draw (wait.south) -- (72.5,12);
  \draw[cbarrow] (68.5,12) -- (action.north);
  \draw[cbarrow,dashed] (action.east) -- (144,0) -- (144,138)
    -- node[above,font=\fontsize{9bp}{11bp}\selectfont] {未停止时更新} (105.5,138) -- (105.5,132);
\end{tikzpicture}

  \bicaption{C-BFC 的整体计算流程}{Overall computation flow of C-BFC}
  \label{rq3:fig:controller}
\end{figure}

这种分工与问题的信息结构相对应。学习模型负责从语义特征中估计任务效果，概率递推负责处理数据是否到达的不确定性，有限前瞻负责选择接下来的通信动作。控制器无需以神经网络重新学习已知的分组转移规则，也不需要在每次重传时重新执行视觉语言模型推理。

\subsection{分类一致的语义效用估计}
\label{rq3:subsec:utility}

\paragraph{集合特征与效用网络}
对于每个非空接收集合 $S$，构造特征向量 $\bm f_k(S)\in\mathbb{R}^{29}$。特征包括融合后的五类信念、最大置信度、类别间隔、熵、证据数量与长度、发送方置信度、局部观测质量、证据一致程度和语义粒度比例。由于此处预测的是“集合 $S$ 已经到达”时的任务效用，送达可靠性按确定接收处理，传输时延不进入语义效用计算；链路的不确定性由接收状态滤波单独处理。

设训练特征的均值和标准差分别为 $\bm\mu$ 与 $\bm s$，先按维度进行标准化：
\begin{equation}
  \widetilde{\bm f}_k(S)=\frac{\bm f_k(S)-\bm\mu}{\max\{\bm s,0.1\}}.
  \label{rq3:eq:feature-normalization}
\end{equation}
效用网络使用两层隐藏表示，结构为
\begin{equation}
  29\longrightarrow128\longrightarrow96\longrightarrow1.
  \label{rq3:eq:network-architecture}
\end{equation}
隐藏层采用 SiLU 激活，训练时使用概率为 0.1 的 Dropout，输出通过 Sigmoid 得到原始效用 $\hat q_\theta(S)$。模型共包含 16,321 个可训练参数。

训练数据由训练任务中所有非空接收集合构成，监督目标为式~\eqref{rq3:eq:true-utility}。设集合样本数量为 $n_{\mathcal D}$，优化目标为二元交叉熵：
\begin{equation}
  \mathcal L(\theta)=-\frac{1}{n_{\mathcal D}}
  \sum_{(k,S)\in\mathcal D}
  \left[u_k(S)\ln\hat q_\theta(S)
  +(1-u_k(S))\ln\bigl(1-\hat q_\theta(S)\bigr)\right].
  \label{rq3:eq:utility-loss}
\end{equation}
训练完成后，通过验证集损失选择模型参数；推理阶段不再使用参考类别。

\paragraph{同类别集合的一致性约束}
独立学习各集合的效用可能破坏实际任务的结构关系。对于固定任务 $k$ 的任意两个非空集合 $S$、$S'$，若 $c_k(S)=c_k(S')$，则 $u_k(S)=u_k(S')$。即使两个集合的置信度或证据数量不同，只要最终输出类别相同，它们相对于同一参考类别的正确性便相同。

基于这一关系，将具有相同预测类别的非空集合划入同一组：
\begin{equation}
  \mathcal G_{k,c}=\{S\in\states_k:S\ne\varnothing,\ c_k(S)=c\}.
  \label{rq3:eq:class-groups}
\end{equation}
在尽量保留原始估计的条件下，求解
\begin{equation}
  \begin{aligned}
    \min_{\{q_k(S)\}}\quad &\sum_{S\ne\varnothing}
       \bigl(q_k(S)-\hat q_\theta(S)\bigr)^2\\
    \mathrm{s.t.}\quad &q_k(S)=q_k(S'),\quad S,S'\in\mathcal G_{k,c}.
  \end{aligned}
  \label{rq3:eq:consistency-projection}
\end{equation}
空集合另外固定为 $q_k(\varnothing)=0$。

\begin{proposition}[一致性投影的闭式解与误差性质]
\label{rq3:prop:projection}
对于任一非空类别组，式~\eqref{rq3:eq:consistency-projection} 的唯一解为组内原始效用的均值：
\begin{equation}
  q_k(S)=\frac{1}{|\mathcal G_{k,c_k(S)}|}
    \sum_{S'\in\mathcal G_{k,c_k(S)}}\hat q_\theta(S').
  \label{rq3:eq:projection-solution}
\end{equation}
若 $\bm u_k$ 是各非空集合的真实正确性向量，则该投影满足
\begin{equation}
  \|\bm q_k-\bm u_k\|_2^2
  \le\|\widehat{\bm q}_k-\bm u_k\|_2^2.
  \label{rq3:eq:projection-error}
\end{equation}
\end{proposition}

\begin{proof}
对类别组 $\mathcal G_{k,c}$，一致性约束要求所有分量取同一数值 $z_c$。组内目标函数为 $\sum_{S\in\mathcal G_{k,c}}(z_c-\hat q_\theta(S))^2$。该函数严格凸，令其导数为零即得到式~\eqref{rq3:eq:projection-solution}。又因为同组集合输出相同类别，其真实效用为同一常数 $u_c$，所以
\begin{equation}
  \sum_{S\in\mathcal G_{k,c}}(\hat q_\theta(S)-u_c)^2
  =\sum_{S\in\mathcal G_{k,c}}(\hat q_\theta(S)-z_c)^2
   +|\mathcal G_{k,c}|(z_c-u_c)^2.
  \label{rq3:eq:projection-decomposition}
\end{equation}
第一项非负，对各组求和即可得到式~\eqref{rq3:eq:projection-error}。
\end{proof}

命题~\ref{rq3:prop:projection} 给出了效用估计在平方误差意义下的性质，实际通信性能还受到能量、时限和状态估计的共同影响。投影使控制器对相同分类结果采用相同效用，避免仅因集合特征变化而继续付出通信代价。$q_k(S)$ 是集合效用分数，不要求不同类别的分数之和为一；严格归一化的概率分布由 $b_{k,t}$ 表示。

\begin{algorithm}[H]
  \bicaption{语义效用模型训练与任务效用表构造}{Semantic utility learning and task utility table construction}
  \label{rq3:alg:utility}
  \small
  \begin{algorithmic}[1]
    \Require 训练与验证语义任务；固定候选生成与融合规则
    \Ensure 效用网络参数 $\theta$；任务 $k$ 的效用表 $q_k$
    \For{每个训练任务及其非空接收集合 $S$}
      \State 提取 $\bm f_k(S)$，按式~\eqref{rq3:eq:true-utility} 计算监督目标
    \EndFor
    \State 计算训练特征统计量，按式~\eqref{rq3:eq:utility-loss} 优化网络
    \State 选择验证交叉熵最低的网络参数并冻结
    \For{待执行任务的每个非空集合 $S$}
      \State 推理 $\hat q_\theta(S)$，计算接收端预测类别 $c_k(S)$
    \EndFor
    \State 按预测类别分组，利用式~\eqref{rq3:eq:projection-solution} 计算 $q_k(S)$
    \State 令 $q_k(\varnothing)=0$
    \State \Return $\theta$ 与 $q_k$
  \end{algorithmic}
\end{algorithm}

\subsection{基于累计确认的接收状态滤波}
\label{rq3:subsec:filter}

在不考虑新反馈到达时，一次发送候选 $m$ 会以概率 $p_{k,m}$ 将当前集合 $S$ 转移为 $S\cup\{m\}$，否则保持原集合。定义信念转移算子
\begin{equation}
  (\mathcal T_m b)(S')=(1-p_{k,m})b(S')
   +p_{k,m}\sum_{S:S\cup\{m\}=S'}b(S).
  \label{rq3:eq:belief-transition}
\end{equation}
该式也适用于重复发送。如果集合原本包含 $m$，则成功或失败均不会改变集合，因此不存在同一证据被重复计入语义信息的问题。

设当前收到的最新累计确认对应第 $j$ 次尝试，其位图为 $S_j^{\mathrm{ack}}$。控制器先将该次尝试结束时的信念恢复为点质量 $\delta_{S_j^{\mathrm{ack}}}$，再重放第 $j+1$ 次到当前第 $n$ 次发送的状态转移：
\begin{equation}
  b_{k,t}=\mathcal T_{m_n}\cdots\mathcal T_{m_{j+1}}
          \delta_{S_j^{\mathrm{ack}}}.
  \label{rq3:eq:ack-replay}
\end{equation}
若尚未收到任何确认，则从空集合点质量出发重放全部发送记录。由于累计位图携带尝试序号，晚到的旧确认不会把系统恢复到过时状态。

\begin{proposition}[接收信念的有效性]
\label{rq3:prop:belief}
对于任意非负、归一化信念 $b$，式~\eqref{rq3:eq:belief-transition} 的结果仍为非负、归一化分布。在第~\ref{rq3:subsec:channel}、\ref{rq3:subsec:feedback} 节的独立性假设下，式~\eqref{rq3:eq:ack-replay} 给出已知累计快照与后续发送历史对应的接收状态条件分布。
\end{proposition}

\begin{proof}
式~\eqref{rq3:eq:belief-transition} 中各项均非负。对 $S'$ 求和时，每个原状态 $S$ 的质量被分配到“不改变集合”和“加入 $m$”两个结果，其权重之和为 $(1-p_{k,m})+p_{k,m}=1$。收到累计确认后，第 $j$ 次尝试结束时的接收集合已经确定；后续尝试的结果在给定链路概率时独立，逐次应用状态转移即可得到当前分布。反馈延迟与丢失独立于数据结果，因此不需要根据未到达的反馈内容进行额外更新。
\end{proof}

例如，在尚无确认且同一证据的成功概率固定为 $p$ 时，连续 $r$ 次尝试后，该证据至少已经送达一次的概率为 $1-(1-p)^r$。控制器可以据此识别重传收益的递减过程，而只统计已确认消息的方法在反馈持续缺失时仍会把其接收概率视为零。

根据接收信念和一致性效用表，当前估计任务效用为
\begin{equation}
  U_k(b)=\sum_{S\in\states_k}b(S)q_k(S).
  \label{rq3:eq:expected-utility}
\end{equation}
后续发送与等待决策均使用这一效用定义。

\subsection{跨任务能量压力与发送价格}
\label{rq3:subsec:queue}

每个节点维护一个虚拟能量队列 $Z_i(k)$，记录相对于平均任务额度的累积超支压力：
\begin{equation}
  Z_i(k+1)=\max\left\{0,Z_i(k)+e_i(k)-\frac{E_i^{(0)}}{N_{\mathrm{task}}}\right\},
  \qquad Z_i(1)=0.
  \label{rq3:eq:virtual-queue}
\end{equation}
节点当前任务的支出高于平均额度时，队列上升；支出较低时，队列逐步回落。因此，即使两个节点的当前链路条件相近，其历史能量压力也可以使控制器倾向于选择不同发送来源。

由于节点剩余电量和待执行任务数持续变化，以当前剩余资源构造平均可用额度
\begin{equation}
  A_{k,t}=\frac{\sum_{i\in\mathcal{N}}E_i(k,t)}{N_{\mathrm{task}}-k+1}.
  \label{rq3:eq:remaining-allowance}
\end{equation}
一次候选发送的资源价格定义为
\begin{equation}
  \pi_{k,m}=\frac{
      \left(1+\dfrac{Z_{s(m)}(k)}{E_{s(m)}^{(0)}}\right)e_{k,m}^{\mathrm d}
      +e_{k,m}^{\mathrm f}}
      {\max\{A_{k,t},\varepsilon_E\}},
  \label{rq3:eq:resource-price}
\end{equation}
其中 $\varepsilon_E$ 为防止分母过小的正数。式~\eqref{rq3:eq:resource-price} 同时反映单次数据与反馈开销、发送节点的历史压力以及剩余任务的平均资源水平。反馈能量进入总通信成本，但不进入感知节点的电量约束。

队列压力与硬电量约束承担不同作用。前者影响动作收益排序，后者直接保证不透支节点电量。各节点独立维护剩余能量，即使系统总余量较大，若某一候选发送节点的剩余能量不足，该动作仍会被式~\eqref{rq3:eq:action-mask} 排除。

\subsection{发送前瞻与反馈等待的联合决策}
\label{rq3:subsec:decision}

为简化记号，本节省略任务与时刻下标，令剩余任务时间为 $h=D_k-t$。发送候选 $m$ 的即时净收益为
\begin{equation}
  g_m(b;h)=U_k(\mathcal T_m b)-U_k(b)
     -\eta_E\pi_{k,m}-\eta_T\frac{\tau_{k,m}}{D_k},
  \label{rq3:eq:send-gain}
\end{equation}
其中 $\eta_E$ 和 $\eta_T$ 分别控制能量价格与显式时间代价。若动作不满足当前时限或节点电量约束，其收益置为负无穷。

\paragraph{两步发送前瞻}
单步收益可能忽略证据之间的组合关系。例如，一条证据单独到达时不能改变最终类别，但它与另一条证据组合后可能产生更好的判断。为保留这种有限组合效应，控制器计算
\begin{equation}
  V_m^{\mathrm{send}}(b;h)=g_m(b;h)
    +\max\left\{0,\max_{m'\in\mathcal F(m)}
      g_{m'}(\mathcal T_m b;h-\tau_{k,m})\right\},
  \label{rq3:eq:two-step-lookahead}
\end{equation}
其中 $\mathcal F(m)$ 按执行 $m$ 后的剩余时间及对应节点电量构造。短前瞻内保持当前能量价格不变，只展开两次数据发送的概率转移；中途反馈的实际到达在下一次重新决策时处理。每次真实发送前均检查尝试次数上限和动作约束。若多个动作的前瞻收益相同，优先选择即时收益较高的动作，再按固定候选次序确定选择。

\paragraph{等待反馈的信息价值}
等待不会直接增加已经送达的证据，但可能揭示当前接收集合，从而改变后续动作。控制器针对最新一次尚未确认的发送快照估计等待价值。设该反馈完成串行发送的时间为 $v_j$，当前反馈年龄为 $a=t-v_j$，额外时延的分布函数为 $F(\cdot)$。在尚未收到这一反馈的条件下，未来等待 $w$ 时间内收到它的概率为
\begin{equation}
  \rho(w)=\frac{(1-\epsilon_{\mathrm f})[F(a+w)-F(a)]}
                {1-(1-\epsilon_{\mathrm f})F(a)}.
  \label{rq3:eq:conditional-feedback}
\end{equation}
分母对应当前仍未收到该反馈的概率，既包含尚未到达，也包含已经丢失两种情况。当反馈已经不可能在未来到达时，继续等待该快照不会产生新的信息价值。

令
\begin{equation}
  G(b;h)=\max\left\{0,\max_m g_m(b;h)\right\}
  \label{rq3:eq:one-step-value}
\end{equation}
表示在给定剩余时间内再执行至多一次发送的最优净收益。若等待后收到最新累计位图，则当前接收集合成为已知；若该反馈仍未到达，则维持现有接收信念。据此，等待收益近似为
\begin{equation}
  \begin{split}
    V^{\mathrm{wait}}(w;b,h)
      ={}&\rho(w)\sum_{S\in\states_k}b(S)G(\delta_S;h-w)\\
         &+[1-\rho(w)]G(b;h-w)-\eta_T\frac{w}{D_k}.
  \end{split}
  \label{rq3:eq:waiting-value}
\end{equation}
该式通过剩余时间 $h-w$ 体现等待的机会成本。即使 $\eta_T=0$，过长等待仍可能使后续发送不再可行，因此等待价值始终受到任务窗口的限制。

等待候选由固定短等待、到反馈平均额外时延的等待和到该分布上界的等待组成：
\begin{equation}
  \mathcal W_t=\left\{w_0,\ \max(w_0,\bar d_{\mathrm f}-a),\
     \max(w_0,1.5\bar d_{\mathrm f}-a)\right\}\cap(0,h).
  \label{rq3:eq:wait-candidates}
\end{equation}
候选仅依据可见发送记录和反馈时延模型生成。实际滤波使用所有已到达的累计确认，而式~\eqref{rq3:eq:waiting-value} 只评价最新未确认快照，忽略其他待到达位图可能提供的部分信息，以限制计算开销。

\paragraph{动作执行与重新决策}
控制器比较所有发送前瞻值、等待值和停止收益零。若最大正收益来自发送，则执行对应候选；若来自等待，则推进相应时间；若不存在正收益，则停止任务。任务结束后按式~\eqref{rq3:eq:virtual-queue} 更新节点队列。完整过程见算法~\ref{rq3:alg:cbfc}。该方法属于基于信念状态的有限前瞻控制，不对全部未来反馈分支或完整任务周期进行精确搜索。

\begin{algorithm}[t]
  \bicaption{C-BFC 在线语义交付控制}{Online semantic delivery control of C-BFC}
  \label{rq3:alg:cbfc}
  \small
  \begin{algorithmic}[1]
    \Require 候选与链路预测；效用模型 $\theta$；初始节点能量；$\eta_E,\eta_T$
    \Ensure 每个任务的发送、等待与停止动作序列
    \State 初始化 $E_i(1)=E_i^{(0)}$，$Z_i(1)=0$
    \For{任务 $k=1,\ldots,N_{\mathrm{task}}$}
      \State 根据算法~\ref{rq3:alg:utility} 构造 $q_k$，令 $t=0$，$b=\delta_{\varnothing}$
      \While{任务尚未停止且 $t<D_k$}
        \State 读取已到达的累计确认与可见发送历史
        \State 按式~\eqref{rq3:eq:ack-replay} 更新接收信念 $b$
        \State 构造可行动作，并按式~\eqref{rq3:eq:resource-price} 计算能量价格
        \If{不存在可行发送动作}
          \State 停止当前任务
        \Else
          \State 按式~\eqref{rq3:eq:two-step-lookahead} 评价各发送动作
          \State 按式~\eqref{rq3:eq:conditional-feedback}--\eqref{rq3:eq:waiting-value} 评价等待候选
          \If{发送与等待的最大收益不大于零}
            \State 停止当前任务
          \Else
            \State 执行收益最大的发送或等待动作
            \State 更新当前时间、发送能量与可见历史
          \EndIf
        \EndIf
      \EndWhile
      \State 汇总节点发送支出，按式~\eqref{rq3:eq:virtual-queue} 更新 $Z_i(k+1)$
    \EndFor
  \end{algorithmic}
\end{algorithm}

\subsection{计算复杂度与适用范围}
\label{rq3:subsec:complexity}

设候选数量为 $K$，效用网络参数量为 $P_\theta$。每个任务需对最多 $2^K$ 个接收集合提取特征并计算效用，该部分只执行一次，可通过批量推理完成。任务内一次信念转移和期望效用计算的开销为 $O(2^K)$；遍历候选并评价两步发送组合的开销约为 $O(K^2 2^K)$。

等待评价需要枚举可能反馈揭示的接收集合，再对各集合评价下一次发送。在直接实现下，其上界为 $O(|\mathcal W_t|K4^K)$。因此，单次在线动作评价的主要复杂度为
\begin{equation}
  O\left(K^2 2^K+|\mathcal W_t|K4^K\right).
  \label{rq3:eq:complexity}
\end{equation}
存储主要包括网络参数、集合特征与效用表、当前接收信念和发送历史。本章取 $K\le4$，因此接收集合不超过十六个，等待候选不超过三个，适合采用显式枚举。

候选规模受限是该实现的重要条件。若大量节点同时产生可发送证据，接收集合规模将指数增长，需要先进行候选压缩，或进一步采用粒子近似与分层控制。此外，本章根据分类结果构造一致性组，其误差性质对应离散分类正确性；对于生成式回答、连续估计或定位误差，应依据相应任务损失重新定义效用约束。

\section{实验设置与结果分析}
\label{rq3:sec:experiments}

\subsection{任务数据与实验流程}
\label{rq3:subsec:experiment-data}

实验采用基于 VQAv2/COCO 图像构造的冻结视觉语义轨迹\citep{coco,vqa-v2}。每个任务由一幅图像及其查询构成，通过局部裁剪为九个感知节点提供不同观测，接收中心不直接访问原图。离线视觉语言模型生成局部证据和全图参考标签，标签空间为清晰、基本清晰、不确定、基本遮挡和遮挡五个等级。因此，本章的任务正确性表示证据可见性分类与全图模型参考标签一致，评价对象是给定冻结语义内容后的通信效果。

候选证据由前期选择性语义通信模块 EviComRL 生成。该模块采用包含结构化候选特征与历史上下文的 Actor-Critic 策略，并以 PPO 进行训练\citep{ppo}。本章使用其在本实验环境训练后固定的检查点，在预设 token 预算下生成候选序列，经固定顺序合并、去重后保留至多四条发往中心的证据。所有比较方法使用相同候选内容与次序。整个反馈控制实验不重新调用视觉语言模型，也不联合微调候选选择模块。

语义效用学习使用训练分区中的 600 个任务，验证分区包含 595 个任务，测试分区包含 611 个任务。图像在不同分区之间不重叠。每个训练任务枚举所有非空接收集合，形成式~\eqref{rq3:eq:utility-loss} 的监督样本；同一图像产生的不同集合仍属于同一个原始任务。效用网络采用 AdamW 优化，学习率为 $8\times10^{-4}$，权重衰减为 0.02，批量大小为 256，最多训练 60 轮，并依据验证交叉熵选择参数。

控制参数的选择在四种验证条件下进行，分别为参考条件、12 ms 反馈延迟、60\% 反馈丢失率和 12 ms 任务时限。对于每种方法及训练种子，首先比较四种条件的平均正确完成率，再以更低通信能耗处理并列情况，最终保留一组参数用于全部测试条件。C-BFC 的三个训练种子均选择 $\eta_E=0.24$、$\eta_T=0$。时间约束仍通过发送可行性、剩余规划时域和等待机会成本影响动作选择。

包含学习模型的方法使用三个训练种子与三个网络种子的九种组合，固定协议使用三个网络种子。每次测试按相同任务次序处理 611 个任务，节点能量在该次任务流内连续扣减。不同方法对同一任务、同一候选和同一次尝试使用相同随机条件，以减少抽样差异对方法比较的影响。图中误差线表示重复实验的标准差。重复实验改变模型初始化和网络随机条件，视觉测试样本保持固定。

\subsection{通信参数与评价指标}
\label{rq3:subsec:experiment-settings}

表~\ref{rq3:tab:parameters} 给出了主要参数。节点到中心的距离在 100--260 m 范围内平滑变化，可用带宽系数在标称带宽的 0.25--1.0 之间变化，分别模拟任务流中传播条件与共享无线资源的波动。速率预测按式~\eqref{rq3:eq:rate} 计算，单次分组成功按式~\eqref{rq3:eq:success-probability} 抽样。每个任务最多执行九次数据发送，不同候选可能来自同一发送节点，因此能量约束始终按节点汇总。

\begin{table}[t]
  \centering\small
  \bicaption{语义交付实验的主要通信与控制参数}{Main communication and control parameters for semantic delivery}
  \label{rq3:tab:parameters}
  \begin{tabularx}{\linewidth}{@{}Yl@{}}
    \toprule
    参数 & 参考值或设置\\
    \midrule
    感知节点数／接收中心数 & 9／1\\
    每任务候选上限 $K_{\max}$ & 4\\
    标称带宽 & 20 kHz\\
    数据／反馈发射功率 & 0.2 W／0.1 W\\
    参考信道增益 $\rho_0$ & $10^{-4}$\\
    噪声功率谱密度 $N_0$ & $4\times10^{-20}$ W/Hz\\
    每 token 字节数 $b$ & 4\\
    数据头部／有效载荷上限 & 64 B／1200 B\\
    反馈帧长度 & 24 B\\
    发送时长裕量 $\kappa$ & 1.6\\
    单次发送时长上限 $\tau_{\max}$ & 8 ms\\
    任务时限 $D_k$ & 25 ms\\
    反馈平均额外时延 $\bar d_{\mathrm f}$ & 6 ms\\
    反馈丢失率 $\epsilon_{\mathrm f}$ & 25\%\\
    速率波动强度 $\sigma_r$ & 0.55\\
    每节点初始发送能量 $E_i^{(0)}$ & 0.15 J\\
    数据尝试上限／基础等待间隔 & 9 次／2 ms\\
    C-BFC 能量／时间权重 & 0.24／0\\
    \bottomrule
  \end{tabularx}
\end{table}

实验使用四项指标共同评价通信效果。

\paragraph{按时正确完成率}
任务必须在截止时间内至少接收一条证据，且最终类别与参考标签一致，才计为正确完成。由于环境仅执行可在截止期内完整结束的数据动作，该指标为
\begin{equation}
  R_{\mathrm c}=\frac{1}{N_{\mathrm{task}}}\sum_{k=1}^{N_{\mathrm{task}}} u_k(S_k^{\mathrm{end}}).
  \label{rq3:eq:metric-correct}
\end{equation}
分母包含因能量不足或策略停止而未获得服务的任务。

\paragraph{单位正确任务通信能耗}
将整个任务流的数据与反馈发送能量除以正确完成的任务数：
\begin{equation}
  E_{\mathrm c}=\frac{\sum_{k=1}^{N_{\mathrm{task}}} E_k^{\mathrm{comm}}}
                         {\sum_{k=1}^{N_{\mathrm{task}}} u_k(S_k^{\mathrm{end}})}.
  \label{rq3:eq:metric-energy}
\end{equation}
该指标以 mJ 表示，可同时反映通信支出与正确交付效果。先对每次回放计算式~\eqref{rq3:eq:metric-energy}，再跨重复实验与条件求平均。

\paragraph{重复发送比例}
若某次发送开始前，同一证据已成功到达接收端，则该次尝试属于重复发送。令 $d_{k,a}$ 为相应指示变量，重复比例为
\begin{equation}
  R_{\mathrm{dup}}=\frac{\sum_k\sum_{a\in\mathcal T_k}d_{k,a}}
                         {\sum_k|\mathcal T_k|}.
  \label{rq3:eq:metric-duplicate}
\end{equation}
这一指标依据实际接收日志计算，仅用于评价控制器未能避免的冗余。已尝试但此前从未成功的补发不计为重复发送。

\paragraph{平均通信完成时延}
记 $\mathcal K_+=\{k:|\mathcal T_k|>0\}$ 为实际发起数据发送的任务集合，则
\begin{equation}
  T_{\mathrm{comm}}=\frac{1}{|\mathcal K_+|}\sum_{k\in\mathcal K_+}T_k.
  \label{rq3:eq:metric-latency}
\end{equation}
该指标包含任务中的发送与等待时间，但不把直接跳过的任务以零时延计入平均。较低时延可能来自及时停止，也可能来自发送较少，因此需要结合正确完成率解释。

\subsection{对比方法与实验分组}
\label{rq3:subsec:baselines}

对比方法包括固定冗余、停等 ARQ、选择重传、语义阈值 ARQ、队列贪心，以及两个学习型控制器。固定冗余按照验证选择的候选数量和副本数发送，不依赖反馈触发后续动作。停等 ARQ 等待累计确认或超时后继续发送；选择重传允许有限窗口内的未确认候选独立补发。语义阈值 ARQ 在累计确认对应的语义效用达到阈值时结束任务。队列贪心使用接收信念与能量价格，逐步选择即时净收益最高的发送动作。

SD-SAC 对照保留文献中的双 Q 平均、Q 值裁剪与熵变化惩罚\citep{zhou2024sdsac}。TP-DQN 参考两阶段码率控制的思路，先使用语义性能预测模型给出任务效用，再训练离散动作价值网络\citep{wu2026twophase}，其价值更新采用 Double DQN 形式。两者均适配为本章的反馈动作空间：四个候选发送动作、三种等待动作和停止动作。策略输入包含候选属性、接收状态概率、语义效用与资源条件，共 89 维，使用 $89\rightarrow128\rightarrow64\rightarrow8$ 的网络结构。

两个学习型控制器都在当前环境重新训练，每个训练种子执行 80 个由 600 个任务构成的任务流，并在十二种物理与反馈设置之间变化。训练每四次交互执行一次更新，批量大小为 128，学习率为 $3\times10^{-4}$。SD-SAC 使用训练任务的正确完成奖励，TP-DQN 使用学习得到的语义性能评价；二者均加入效用势函数差作为训练信号。最终测试使用验证选中的在线训练参数。这里比较的是相同语义交付任务中的控制机制，原研究中的文本编解码网络不属于本实验环境。

实验按五组条件展开：速率波动强度取 0.20、0.35、0.55、0.80、1.05；反馈平均额外时延取 0、3、6、9、12 ms；任务时限取 10、15、25、35、45 ms；每节点发送能量取 0.06、0.10、0.15、0.24、0.36 J；组件实验在 0\%、15\%、25\%、45\%、60\% 五档反馈丢失率下进行。每组只改变对应参数，其他参数保持参考值。前四组去除重复参考点后共有十七个主比较条件，连同组件实验共二十一个不同条件，总计完成 1,089 次完整任务流回放。

主图展示固定冗余、语义阈值 ARQ、队列贪心、SD-SAC、TP-DQN 和 C-BFC，性能表同时列出停等 ARQ 与选择重传结果。含语义估计的对照使用相同训练种子的原始效用模型，C-BFC 在此基础上执行一致性投影。各方法的阈值、重传规则或资源权重均在验证分区选择后固定。

\subsection{总体性能比较}
\label{rq3:subsec:overall-results}

表~\ref{rq3:tab:main-results} 给出了十七个主条件的等权平均结果。C-BFC 的正确完成率为 49.84\%，单位正确任务通信能耗为 3.615 mJ。与固定冗余和语义阈值 ARQ 相比，C-BFC 分别提高正确完成率 9.35 和 6.45 个百分点。固定发送副本能够提高部分证据的交付机会，但无法根据已发生的传输结果及时调整；语义阈值 ARQ 能够利用已确认集合的效用停止发送，但在反馈未返回时仍主要受超时规则支配。

\begin{table}[t]
  \centering\small
  \setlength{\tabcolsep}{3.5pt}
  \bicaption{十七个主比较条件下的平均性能}{Average performance across the seventeen main evaluation conditions}
  \label{rq3:tab:main-results}
  \begin{tabular*}{\linewidth}{@{\extracolsep{\fill}}lrrrr@{}}
    \toprule
    方法 & \shortstack{正确完成率\\(\%)} & \shortstack{单位正确任务\\能耗 (mJ)} & \shortstack{重复发送\\(\%)} & \shortstack{通信时延\\(ms)}\\
    \midrule
    固定冗余 & 40.49 & 4.941 & 29.35 & 11.32\\
    停等 ARQ & 43.65 & 3.678 & 7.49 & 16.69\\
    选择重传 & 43.57 & 3.665 & 7.16 & 16.77\\
    语义阈值 ARQ & 43.39 & 4.090 & 6.71 & 17.63\\
    队列贪心 & 47.26 & 4.564 & 25.76 & 10.96\\
    SD-SAC & 48.97 & 4.334 & 15.67 & 17.76\\
    TP-DQN & 49.18 & 4.219 & 19.11 & 17.13\\
    \midrule
    \textbf{C-BFC} & \textbf{49.84} & \textbf{3.615} & \textbf{17.14} & \textbf{16.47}\\
    \bottomrule
  \end{tabular*}
  \par\vspace{4pt}{\footnotesize 各条件等权汇总，共享的参考条件只计入一次。加粗行表示本章方法。}
\end{table}


队列贪心的正确完成率为 47.26\%，说明显式接收信念和语义收益能够改善简单重传规则的任务效果。但是，其单位正确任务能耗为 4.564 mJ，重复发送比例为 25.76\%。这一结果与其仅评价当前发送收益、缺少主动等待相符：当接收状态尚不确定时，连续执行即时收益为正的动作，可能消耗本可通过等待确认而节省的能量。

相较于 SD-SAC 和 TP-DQN，C-BFC 的正确完成率分别提高 0.87 和 0.66 个百分点，单位正确任务能耗分别降低 16.59\% 和 14.32\%，通信完成时延分别降低 7.29\% 和 3.87\%。在十七个主条件中，C-BFC 相对于两者的正确完成率分别在十二个和十四个条件下更高，单位正确任务能耗在全部十七个条件下更低。其平均正确率优势较小，主要收益表现为在相近或更高正确交付水平下减少通信支出。

\begin{table}[t]
  \centering\small
  \setlength{\tabcolsep}{3.5pt}
  \bicaption{参考条件下的性能比较}{Performance comparison under the reference condition}
  \label{rq3:tab:reference-results}
  \begin{tabular*}{\linewidth}{@{\extracolsep{\fill}}lrrrr@{}}
    \toprule
    方法 & \shortstack{正确完成率\\(\%)} & \shortstack{单位正确任务\\能耗 (mJ)} & \shortstack{重复发送\\(\%)} & \shortstack{通信时延\\(ms)}\\
    \midrule
    固定冗余 & 41.68 & 5.023 & 29.86 & 11.79\\
    停等 ARQ & 45.66 & 3.627 & 7.72 & 17.46\\
    选择重传 & 45.50 & 3.615 & 7.40 & 17.46\\
    语义阈值 ARQ & 46.32 & 3.948 & 7.07 & 18.16\\
    队列贪心 & 47.66 & 4.706 & 27.20 & 11.36\\
    SD-SAC & 49.95 & 4.356 & 15.66 & 18.24\\
    TP-DQN & 49.86 & 4.329 & 21.49 & 17.74\\
    \midrule
    \textbf{C-BFC} & \textbf{51.08} & \textbf{3.639} & \textbf{17.40} & \textbf{17.56}\\
    \bottomrule
  \end{tabular*}
  \par\vspace{4pt}{\footnotesize 参考条件为 25 ms 任务时限、6 ms 反馈延迟、25\% 反馈丢失率及每节点 0.15 J 发送能量。}
\end{table}


参考条件下的结果见表~\ref{rq3:tab:reference-results}。C-BFC 的正确完成率为 51.08\%，高于两个学习型对照约 1.1 个百分点，单位正确任务能耗由约 4.3 mJ 降至 3.639 mJ。停等与选择重传具有较低重复比例，说明保守确认规则有助于避免对已接收证据再次发送；与此同时，其正确完成率约为 45.5\%，表明较少重复并不必然对应更好的任务完成效果。因而，本章将正确完成率与单位正确任务能耗作为主要结果，并通过重复比例和时延解释其形成过程。

\FloatBarrier
\subsection{信道波动对交付效果的影响}
\label{rq3:subsec:fading-results}

图~\ref{rq3:fig:fading} 比较不同速率波动强度下的性能。在预测平均速率不变时，较大的 $\sigma_r$ 使单次发送更容易无法在预留时长内完成，因而多数方法的正确完成率下降，单位正确任务能耗上升。C-BFC 的正确完成率由 $\sigma_r=0.20$ 时的 57.26\% 下降至 $\sigma_r=1.05$ 时的 41.32\%，说明有限重传与反馈控制仍然受数据链路可靠性的约束。

\begin{figure}[!htbp]
  \centering
  \resultfigure{01_fading.pdf}{不同速率波动强度下的性能比较}
  \bicaption{速率波动强度变化下的性能比较}{Performance under different rate fluctuation levels}
  \par\vspace{2pt}{\footnotesize （a）按时正确完成率；（b）单位正确任务通信能耗；（c）重复发送比例；（d）平均通信完成时延。其余参数取表~\ref{rq3:tab:parameters} 的参考值。}
  \label{rq3:fig:fading}
\end{figure}

信道较稳定时，C-BFC 的节能优势更为明显。在 $\sigma_r=0.20$ 时，其单位正确任务能耗为 2.346 mJ，SD-SAC 和 TP-DQN 分别为 3.464 和 2.831 mJ，而三个方法的正确完成率均在 57\% 左右。此时，多次重传的实际需求较低，对接收状态与分类结果一致性进行判断能够较早识别额外发送的有限收益。

当 $\sigma_r=1.05$ 时，C-BFC 的单位正确任务能耗增至 5.036 mJ，仍低于 SD-SAC 的 5.621 mJ 和 TP-DQN 的 5.610 mJ。与此同时，各方法之间的正确完成率差距缩小。该结果表明，控制策略可以改善有限资源的使用方式，但无法完全抵消低成功概率对语义交付造成的损失。

\FloatBarrier
\subsection{确认反馈延迟的影响}
\label{rq3:subsec:feedback-results}

图~\ref{rq3:fig:feedback} 给出了反馈平均额外时延变化下的结果。额外时延为零时，反馈仍需完成串行发送，因而并非瞬时获得确认。随着额外时延增加，发送端维持接收不确定性的时间变长，等待确认与继续发送之间的矛盾更加明显。

\begin{figure}[!htbp]
  \centering
  \resultfigure{02_feedback.pdf}{不同反馈延迟下的性能比较}
  \bicaption{反馈平均额外时延变化下的性能比较}{Performance under different additional feedback delays}
  \par\vspace{2pt}{\footnotesize 四个子图依次对应正确完成率、单位正确任务能耗、重复发送比例和通信完成时延。}
  \label{rq3:fig:feedback}
\end{figure}
\FloatBarrier

在额外时延为 0 ms 时，C-BFC 的正确完成率为 53.52\%，通信完成时延为 11.36 ms；SD-SAC 与 TP-DQN 的正确完成率分别为 50.95\% 和 52.01\%，通信完成时延分别为 15.60 和 14.42 ms。较及时的确认能够使控制器尽早确定已经到达的证据，并据此停止冗余发送或选择互补内容。

当反馈平均额外时延增至 12 ms，C-BFC 的正确完成率降至 48.92\%，单位正确任务能耗由 3.447 mJ 增至 3.973 mJ，重复发送比例也随之增加。此时部分确认无法在有用的决策窗口内返回，主动等待的作用受到限制。两个学习型对照在同一条件下的正确完成率分别为 48.10\% 和 48.52\%，单位正确任务能耗分别为 4.426 和 4.526 mJ。C-BFC 仍保持较低的单位正确任务能耗，但其相对收益没有随着反馈延迟无限扩大。

\FloatBarrier
\subsection{任务截止时间的影响}
\label{rq3:subsec:deadline-results}

图~\ref{rq3:fig:deadline} 比较不同任务时限下的性能。时限同时决定能够执行的数据尝试数量和可用于等待确认的时间。10 ms 条件下，许多任务只能够容纳一次主要发送，控制器很难通过“发送后等待、再依据反馈补充证据”的过程获得额外收益。C-BFC 与 SD-SAC 的正确完成率分别为 44.48\% 和 44.64\%，差异较小。

\begin{figure}[!htbp]
  \centering
  \resultfigure{03_deadline.pdf}{不同任务截止时间下的性能比较}
  \bicaption{任务截止时间变化下的性能比较}{Performance under different task deadlines}
  \par\vspace{2pt}{\footnotesize 较长任务窗口允许反馈到达后继续调整发送，各方法均受到跨任务节点电量限制。}
  \label{rq3:fig:deadline}
\end{figure}

随着时限增加，C-BFC 的正确完成率提升，同时单位正确任务能耗保持在较低水平。在 45 ms 条件下，其正确完成率达到 54.74\%，较 SD-SAC 的 51.34\% 提高 3.40 个百分点；单位正确任务能耗为 3.539 mJ，较 SD-SAC 的 4.466 mJ 降低 20.76\%。与 TP-DQN 相比，C-BFC 也在正确完成率和单位正确任务能耗上取得改善。

C-BFC 在 45 ms 任务窗口下的平均通信完成时延为 23.39 ms，低于 SD-SAC 的 28.20 ms 和 TP-DQN 的 26.03 ms。算法能够利用返回的确认调整证据发送，并在继续通信的收益不足时提前结束。结合紧时限结果可见，反馈控制的作用取决于任务窗口能否容纳信息返回后的后续动作。

\FloatBarrier
\subsection{节点发送能量的影响}
\label{rq3:subsec:energy-results}

图~\ref{rq3:fig:energy} 将每节点初始发送能量从 0.06 J 增加到 0.36 J。能量额度较低时，当前任务中的支出会较快影响后续任务的可执行范围，正确完成率受到整个任务流资源总量的限制。在每节点 0.06 J 条件下，C-BFC 的正确完成率为 36.32\%，单位正确任务能耗为 2.427 mJ，说明其通过限制低收益发送保持了较高的资源使用效率。

\begin{figure}[!htbp]
  \centering
  \resultfigure{04_energy.pdf}{不同节点发送能量下的性能比较}
  \bicaption{每节点初始发送能量变化下的性能比较}{Performance under different initial node transmission energies}
  \par\vspace{2pt}{\footnotesize 电量在同一次 611 个任务的回放中连续扣减，其他通信条件固定。}
  \label{rq3:fig:energy}
\end{figure}

随着能量增加，各方法能够为更多任务提供通信机会，也能够执行更多补发。C-BFC 在每节点 0.10 J 时的正确完成率为 46.03\%，高于 SD-SAC 的 44.66\%；在参考能量 0.15 J 下进一步达到 51.08\%。这一变化说明，虚拟能量压力与剩余额度能够使同一组控制权重随资源条件改变实际发送行为。

在较高能量额度下，需要同时观察质量与能耗。每节点 0.36 J 时，C-BFC 的正确完成率为 55.74\%，低于 SD-SAC 的 57.57\% 和 TP-DQN 的 56.76\%；对应单位正确任务能耗为 4.608 mJ，低于两者的 6.164 和 6.033 mJ。此时，C-BFC 保留的能量偏好使其放弃了一部分代价较高的交付机会。因而，高能量条件下的结果体现了正确率与通信消耗之间的取舍，不能仅根据节能幅度判断所有性能指标均有改善。

\FloatBarrier
\subsection{关键组件作用分析}
\label{rq3:subsec:ablation}

为分析主要组件的作用，构造三个对照版本：无接收信念版本只使用最新确认位图，不传播未确认尝试的成功概率，也不计算依赖该分布的等待信息价值；无主动等待版本保留发送前瞻，但不将等待作为可主动选择的动作；无一致性约束版本直接使用原始效用预测。各版本在验证分区分别选择资源权重，再在五档反馈丢失率下进行测试。图~\ref{rq3:fig:ablation} 展示逐条件变化，表~\ref{rq3:tab:ablation-results} 给出等权平均结果。

\begin{figure}[!htbp]
  \centering
  \resultfigure{05_ablation.pdf}{不同反馈丢失率下的组件比较}
  \bicaption{不同反馈丢失率下的组件比较}{Component comparison under different feedback loss rates}
  \par\vspace{2pt}{\footnotesize 完整方法与无接收信念、无主动等待、无一致性约束三个版本分别比较，子图指标与前四组一致。}
  \label{rq3:fig:ablation}
\end{figure}

\begin{table}[t]
  \centering\small
  \setlength{\tabcolsep}{3.5pt}
  \bicaption{五档反馈丢失率下的组件比较}{Component comparison at five feedback loss rates}
  \label{rq3:tab:ablation-results}
  \begin{tabular*}{\linewidth}{@{\extracolsep{\fill}}lrrrr@{}}
    \toprule
    方法 & \shortstack{正确完成率\\(\%)} & \shortstack{单位正确任务\\能耗 (mJ)} & \shortstack{重复发送\\(\%)} & \shortstack{通信时延\\(ms)}\\
    \midrule
    无接收信念 & 39.69 & 5.814 & 60.46 & 16.91\\
    无主动等待 & 48.89 & 4.516 & 32.59 & 10.96\\
    无一致性约束 & 51.28 & 4.064 & 16.96 & 18.63\\
    \midrule
    \textbf{C-BFC} & \textbf{50.97} & \textbf{3.674} & \textbf{18.29} & \textbf{17.62}\\
    \bottomrule
  \end{tabular*}
  \par\vspace{4pt}{\footnotesize 各版本分别完成验证参数选择；表中数值为五个反馈丢失条件的等权平均。}
\end{table}


去除接收信念后，平均正确完成率由完整方法的 50.97\% 降至 39.69\%，单位正确任务能耗由 3.674 mJ 增至 5.814 mJ，重复发送比例达到 60.46\%。只依赖已确认位图会使控制器在反馈缺失时持续低估接收端掌握的信息，进而反复向已接收证据分配资源。这一结果支持显式区分数据丢失与反馈缺失的必要性。

去除主动等待后，平均正确完成率为 48.89\%，单位正确任务能耗为 4.516 mJ，重复发送比例为 32.59\%。完整方法通过等待可能到达的确认，减少了状态不确定时的连续发送，但平均通信完成时延也由 10.96 ms 增至 17.62 ms。可见，等待的价值表现为用一部分可用时间换取更合适的后续发送选择，其作用不能仅通过时延这一指标衡量。

去除一致性约束后，平均正确完成率为 51.28\%，略高于完整方法的 50.97\%；单位正确任务能耗为 4.064 mJ，高于完整方法的 3.674 mJ。加入一致性约束使单位正确任务能耗降低约 9.6\%，并伴随 0.31 个百分点的正确率变化。该结果与一致性投影的设计目的相符：抑制不改变分类结果的估计效用增量，主要改善通信支出，而不要求每个条件的正确率都提高。

\FloatBarrier
\subsection{结果讨论}
\label{rq3:subsec:discussion}

综合五组实验，C-BFC 的作用主要体现在接收不确定性、反馈等待与语义冗余共同存在的条件下。接收信念减少对确认缺失的误判，等待价值决定何时值得为信息返回保留时间，一致性效用使资源分配与接收端分类结果保持联系。三个部分共同改善了单位正确任务的通信消耗，但其收益受到链路成功概率、任务窗口和剩余电量的约束。

本章实验也给出了该方法的适用条件。首先，任务语义来自冻结轨迹，能够保证不同控制器面对相同候选与融合规则，但没有覆盖视觉语言模型在线生成内容发生变化的情况。其次，接收集合依赖至多四条候选，状态枚举的开销受到这一规模限制。再次，数据与确认采用独立抽样的发送模型，尚未考虑强相关衰落、反馈冲突或软信息合并。因此，本章结论对应调度层的语义可靠交付，不直接等同于完整无线原型系统的性能。

从资源配置角度看，较短时限首先限制反馈的可用程度，较低能量则限制整个任务流中的服务机会；资源较充足时，控制器仍需在少量额外正确完成与较高通信支出之间选择。本章通过固定验证参数考察这种条件变化，后续可进一步研究面向不同服务质量要求的工作点选择，以适应可靠性与能耗偏好不同的应用。


\section{本章小结}

本章以效用一致性、接收信念与有限前瞻构建 C-BFC，协调不完备反馈下的语义交付与跨任务能量支出。十七个主条件下，平均正确完成率为 49.84\%，单位正确任务能耗为 3.615 mJ，较 SD-SAC 和 TP-DQN 分别降低 16.59\% 和 14.32\%。方法在相近交付水平下减少通信支出，收益仍受时限与能量条件制约。
